\Askhsh{20091}{4}
\begin{ylh}{1}
    \item 1.4. Συντεταγμένες στο Επίπεδο
    \item 2.1. Εξίσωση Ευθείας
    \item 3.1 Ο Κύκλος.
\end{ylh}
Τα σημεία $A(-7, -1)$ και $B(3, -5)$ είναι σημεία ενός κύκλου $C$ κέντρου $K$. Το σημείο $M$ είναι το μέσο της χορδής $AB$ και μία ευθεία $\varepsilon$ διέρχεται από τα σημεία $K$ και $M$.
\begin{alist}
    \item Να βρείτε:
    \begin{rlist}
        \item Τις συντεταγμένες του σημείου $M$. \monades{04}
        \item Την εξίσωση της ευθείας $KM$. \monades{08}
    \end{rlist}
    \item Αν από το κέντρο $K$ του κύκλου διέρχεται η ευθεία $(\delta): x+y = -12$, τότε:
    \begin{rlist}
        \item Να βρείτε τις συντεταγμένες του σημείου $K$. \monades{07}
        \item Να βρείτε την εξίσωση του κύκλου $C$. \monades{06}
    \end{rlist}
\end{alist}